\documentclass{ou}
\begin{document}

\thispagestyle{empty}
\begin{center}
\shortstack{\FGAVO}\vspace{4\baselineskip}
\textbf{РЕФЕРАТ}\\
\setstretch{2}
Информационные системы и~технологии\\
\setstretch{1}
Дополненная и виртуальная реальность\vspace{2\baselineskip}\end{center}
\begin{flushleft}
\doublespacing
\begin{tabularx}{\textwidth}{@{} p{0.4\textwidth} l l @{}}
Руководитель & ст. преподаватель & И. Р. Нигматуллин\\
Студенты     &                   &                  \\
\multirow{-2.27}{\hsize}{\singlespacing Представление трехмерных объектов как формы цифровой информации}
             & ГФ25-01Б, 152540765 & А. В. Гайворонская\\
             & ГФ25-01Б, 152540753 & М. И. Бобылев\\
             & ГФ25-01Б, 152539850 & Е. А. Аксаментова\\
\multirow{-2.27}{\hsize}{\singlespacing Основы 3D-моделирования и~фотограмметрии для~создания виртуальных объектов}
             & ГФ25-01Б, 152537671 & А. Е. Корикова\\
             & ГФ25-01Б, 152536528 & Д. А. Васильева\\
             & ГФ25-01Б, 152540016 & И. Н. Морозова\\         
\multirow{-2.27}{\hsize}{\singlespacing Принципы функционирования технологий дополненной и~виртуальной реальности}
             & ГФ25-01Б, 152537684 & О. И. Бочанова\\
             & ГФ25-01Б, 152539986 & М. К. Гинтер\\
             & ГФ25-01Б, 152541471 & А. О. Безруких\\
             & ГФ25-01Б, 152536796 & А. В. Корюкова\\
\end{tabularx}
\end{flushleft}
\vspace*{\fill}
\begin{center}
\city
\end{center}
\newpage
\thispagestyle{empty}
\setstretch{1}
Продолжение титульного~листа реферата по~теме «Дополненная и~виртуальная реальность»
\begin{flushleft}
\doublespacing
\begin{tabularx}{\textwidth}{@{} p{0.4\textwidth} l l @{}}
Студенты     &                     &         \\
\multirow{-2.27}{\hsize}{\singlespacing Среды и инструменты для~разработки контента в~AR/VR}
             & ГФ25-01Б, 152541294 & Э. Р. Махмудов\\
             & ГФ25-01Б, 152539829 & В. Д. Бобкова\\
\multirow{-2.27}{\hsize}{\singlespacing Применение AR и VR в~гуманитарных исследованиях и образовании}
             & ГФ25-01Б, 152536115 & С. В. Идрисова\\
             & ГФ25-01Б, 152540131 & Е. Кичко\\
             & ГФ25-01Б, 152539526 & К. Н. Аркадьева\\
             & ГФ25-01Б, 152536056 & В. А. Карасова\\
\multirow{-2.27}{\hsize}{\singlespacing Тенденции развития иммерсивных технологий и их влияние на современную культуру}
             & ГФ25-01Б, 152540046 & А. Д. Балцевич\\
             & ГФ25-01Б, 152537659 & В. Е. Зайцева\\
             & ГФ25-01Б, 152536592 & А. Ю. Бирюкова\\   
\end{tabularx}
\end{flushleft}
\newpage
{
\sloppy
\tableofcontents
}
\newpage
\onehalfspacing
\section{Представление трехмерных объектов как формы цифровой информации}
\subsection{Понятие трехмерного объекта как цифровой информации}
Трёхмерный объект в цифровой среде -- это математическое описание структуры, формы и свойств предмета, которое позволяет визуализировать его на экране и взаимодействовать с ним в виртуальном пространстве.

В отличие от 2D-изображений, представляющих плоскость пикселей, 3D-модели основаны на координатах x, y, z, задающих положение точек в~пространстве. На основе этих данных строится геометрическая форма.

3D-объект включает:
\begin{itemize}[label=-, leftmargin=0em, itemindent=3.35em, nosep]
\item геометрию -- форму, контуры, объём;
\item текстуры -- изображения, накладываемые на поверхность;
\item материалы -- свойства поверхности (блеск, прозрачность, отражение);
\item анимационные данные -- кости, скелеты, ключевые кадры;
\item метаданные -- масштаб, авторство, параметры освещения.
\end{itemize}

Такое представление делает 3D-объекты мощным инструментом для~моделирования и анализа.
\subsection{Основные способы представления трехмерных объектов}
Существует несколько подходов к описанию 3D-модели. У каждого есть свои особенности, преимущества и области применения.

Самый распространённый тип -- полигональная сетка. Объект состоит из:
\begin{itemize}[label=-, leftmargin=0em, itemindent=3.35em, nosep]
\item вершин -- точек в пространстве;
\item рёбер -- соединений между точками;
\item полигонов -- чаще всего треугольников и четырёхугольников.
\end{itemize}

Чем больше полигонов -- тем выше детализация. Преимущества:
\begin{enumerate}[leftmargin=0em, itemindent=3.75em, nosep, label=\arabic*)]
\item Высокая скорость рендеринга;
\item Простота обработки;
\item Широко поддерживается всеми графическими движками.
\end{enumerate}

Недостатки:
\begin{enumerate}[leftmargin=0em, itemindent=3.75em, nosep, label=\arabic*)]
\item Сложность представления идеально гладких поверхностей;
\item Высокая нагрузка при большом количестве полигонов.
\end{enumerate}

Пример: в играх 2000-х персонажи имели около 2--5 тысяч полигонов. Сегодня модели для AAA-игр достигают 50--150 тысяч, а иногда и нескольких миллионов.

NURBS -- это математические кривые и поверхности, создаваемые на~основе контрольных точек. Они позволяют получать идеально гладкие формы. Используются в:
\begin{itemize}[label=-, leftmargin=0em, itemindent=3.35em, nosep]
\item промышленном дизайне;
\item автомобилестроении;
\item моделировании корпусов кораблей;
\item создании ювелирных изделий.
\end{itemize}

Пример: большинство деталей автомобилей сначала проектируют в~виде NURBS-поверхностей, а затем превращают в полигональные модели для~визуализации.

Воксели -- это 3D-аналог пикселей, представляющий объект в виде объёма, заполненного элементами одинакового размера. Преимущества:
\begin{enumerate}[leftmargin=0em, itemindent=3.75em, nosep, label=\arabic*)]
\item Подходит для моделирования природных объектов: облаков, тканей, почвы, органов;
\item Хорошо работает в медицинской визуализации (КТ, МРТ);
\item Простая структура.
\end{enumerate}

Недостатки:
\begin{enumerate}[leftmargin=0em, itemindent=3.75em, nosep, label=\arabic*)]
\item Большие объёмы данных;
\item Медленная обработка.
\end{enumerate}

Пример: игра Minecraft работает на вокселях -- каждый блок -- это воксель, который можно изменить.

Параметрические модели бъект задаётся математическими параметрами, например:
\begin{itemize}[label=-, leftmargin=0em, itemindent=3.35em, nosep]
\item радиус;
\item высота;
\item количество сегментов;
\item формулы генерации.
\end{itemize}

Такие модели легко изменять -- достаточно поменять числовые значения.

Пример: в CAD-системах архитектор может удлинить колонну, просто меняя её высоту, не перестраивая вручную всю геометрию.

Процедурная генерация создаёт объекты по алгоритму. Например:
\begin{enumerate}[leftmargin=0em, itemindent=3.75em, nosep, label=\arabic*)]
\item Генерация ландшафтов;
\item Создание городов;
\item Моделирование космических объектов;
\item Дизайн растительности.
\end{enumerate}

Пример: в игре No Man’s Sky практически вся вселенная создана процедурно -- каждое дерево, планета и животное генерируются алгоритмом.
\subsection{Текстуры, материалы и карты}
Геометрии недостаточно -- чтобы 3D-объект выглядел реалистично, используются дополнительные данные. Типы карт:
\begin{itemize}[label=-, leftmargin=0em, itemindent=3.35em, nosep]
\item Diffuse map -- основной цвет;
\item Normal map -- имитация неровностей;
\item Specular map -- отражение света;
\item Displacement map -- реальное смещение поверхности;
\item Ambient occlusion -- мягкие тени в углублениях.
\end{itemize}
\subsection{Форматы хранения 3D-объектов}
Существует множество форматов, каждый из которых подходит под~конкретные задачи. Популярные форматы:
\begin{enumerate}[leftmargin=0em, itemindent=3.75em, nosep, label=\arabic*)]
\item OBJ -- универсальный, простой текстовый формат. Используется в~играх, моделировании;
\item FBX -- поддерживает анимацию, материалы. Используется в кино, анимации, Unity, Unreal Engine;
\item STL -- содержит только геометрию. Используется в 3D-печати;
\item GLTF/GLB -- оптимизирован для веба;
\item STEP -- параметрические данные. Используется в инженерии.	
\end{enumerate}

Пример: Для 3D-печати почти всегда используют формат STL, потому~что~он прост и передаёт только геометрию, необходимую для печати.
\subsection{Области применения трехмерных объектов}
В промышленности и инженерии 3D-модели позволяют конструировать:
\begin{itemize}[label=-, leftmargin=0em, itemindent=3.35em, nosep]
\item двигатели;
\item корпуса машин;
\item детали механизмов.
\end{itemize}

Перед выпуском изделия создают цифровой двойник и тестируют его в~виртуальной среде.

Архитекторы проектируют здания, включая:
\begin{itemize}[label=-, leftmargin=0em, itemindent=3.35em, nosep]
\item конструкции;
\item систему коммуникаций;
\item интерьер.
\end{itemize}

Это ускоряет расчёты и помогает визуализировать проект заказчику.

В медицине 3D-модели используются для:
\begin{itemize}[label=-, leftmargin=0em, itemindent=3.35em, nosep]
\item подготовки к операциям;
\item моделирования органов;
\item создания протезов.
\end{itemize}

КТ и МРТ дают воксельные данные, которые преобразуются в 3D-модели.

Фильмы, игры, VR-приложения основаны на трёхмерных моделях. Пример: в фильмах Marvel актёры снимаются перед зелёным экраном, а все окружение создаётся 3D-художниками.

Примеры реальных 3D-моделей:
\begin{enumerate}[leftmargin=0em, itemindent=3.75em, nosep, label=\arabic*)]
\item 3D-модель скелета человека для обучения студентов-медиков -- позволяет рассматривать кости под любым углом, изучать структуру;
\item Модель здания для строительной компании -- заказчик может «пройтись» по будущей квартире;
\item Модель двигателя самолёта -- инженеры тестируют работу частей, прежде чем производить их.
\item Игровой персонаж -- состоит из сетки, текстур, скелета и анимаций.
\end{enumerate}
\newpage
\section{Основы 3D-моделирования и~фотограмметрии для~создания виртуальных объектов}
Современные подходы к созданию виртуальных объектов основываются на~двух принципиально различных, но взаимодополняющих методологиях -- 3D-моделировании и фотограмметрической реконструкции. Эти технологии формируют основу производственных процессов в таких областях, как разработка игр, кинематограф, виртуальная и дополненная реальность, цифровое сохранение культурного наследия и промышленный дизайн.

3D-моделирование -- процесс искусственного конструирования трёхмерных объектов через манипулирование геометрическими примитивами. Данный подход обеспечивает полный контроль над создаваемой формой, позволяя проектировать оптимизированные для анимации и рендеринга модели с заданными параметрами топологии. Этот метод остаётся основным для создания абстрактных, стилизованных или несуществующих в реальности объектов\cite{gb_3d}.

Фотограмметрия решает задачу автоматизированного переноса реальных объектов в цифровую форму на основе алгоритмического анализа множества перекрывающихся изображений. Метод позволяет воссоздавать трёхмерную геометрию с высокой метрической точностью и~фотореалистичным текстурным покрытием, однако результирующие модели требуют значительной постобработки для их дальнейшего практического использования\cite{sky_fotogrammetriya}.
\subsection{Что такое 3D-моделирование}
3D моделирование -- процесс создания трехмерных цифровых объектов в виртуальном пространстве. В отличие от 2D графики, трехмерные модели обладают глубиной, объемом и могут быть рассмотрены с любого ракурса -- словно реальные физические объекты\cite{sky_3d}.

Базовая концепция 3D моделирования строится на координатной системе, где каждая точка определяется тремя значениями (X, Y, Z), отвечающими за~ширину, высоту и глубину. Это формирует основу для создания виртуальных объектов любой сложности.

Хотя программное обеспечение для 3D-моделирования основано на~сложных математических расчетах, вычисления проводятся автоматически с предоставлением удобного пользовательского интерфейса. Создание трехмерной модели довольно затруднительно и представляет собой своего рода искусство. Для достижения реалистичности необходимо разбираться в~особенностях моделирования и правильно проводить расчеты в течение всего процесса моделирования.
\subsection{Принципы 3D-моделирования}
Создание трехмерных моделей основано на двух основных принципах наглядности и информативности.

Наглядность -- это свойство изображения, состоящее в правильном и~четком представлении о моделируемом объекте. Наглядность достигается путем внешнего оформления трехмерной модели, которое включает в себя цвета, обозначения, форму и размер элементов, текстуру, то есть наглядность предполагает возможность восприятия зрителем форм, размеров и содержания трехмерной модели.

Чем детальнее модель, тем больше она включает элементов с~большими подробностями. Наглядность изображения повышается путем оптимизации данных, то есть скрытия объектов, являющихся несущественными.

Информативность -- свойство трехмерных изображений, зависящее от~количества содержащихся в них пространственных характеристик. Наибольшая информативность трехмерного изображения достигается при~всестороннем представлении внешнего вида, положения в пространстве, размеров и форм всех значимых элементов модели\cite{sushentsov_2023}.

Виды 3D-моделирования:
\begin{itemize}[label=-, leftmargin=0em, itemindent=3.35em, nosep]
\item Полигональное 3D-моделирование;
\item Параметрическое 3D-моделирование.
\end{itemize}

Полигональное моделирование заключается в построении трехмерной фигуры на основе плоской поверхности, которая размечается сеткой. Сетка состоит из линий, называемых ребрами, которые пересекаются в точках, называемых вершинами. Ребра делят поверхность на отдельные полигоны.

На программном уровне осуществляются действия с ребрами, вершинами до тех пор, пока объект не примет нужную форму, при этом происходит смещение полигонов относительно друг друга под различными углами. Число полигонов может достигать огромных значений. По мере его увеличения сетка все сильнее напоминает контуры создаваемого объекта, и он все более становится таким, как задумывалось.

Это аналогично тому, как правильный многоугольник при увеличении числа ребер принимает форму круга.

Также в качестве полигонов могут выступать отдельные двухмерные фигуры, называемые сплайнами. Они могут выглядеть как простые фигуры, отдельные фигуры и линии, так и составные. Вместе они соединяются в одну трехмерную фигуру. Такой способ моделирования уместен, если автор хочет, чтобы зритель увидел элементы, образующие 3D-фигуру.

Полигональное 3D-моделирование встречается чаще всего используется в таких областях как:
\begin{itemize}[label=-, leftmargin=0em, itemindent=3.35em, nosep]
\item наука;
\item архитектура;
\item компьютерные игры;
\item дополненная и виртуальная реальности;
\item 3D-печать;
\item графические элементы для веб-интерфейса (смайлы, кнопки);
\item спецэффекты в фильмах;
\item скульптинг (статуи, скульптуры)\cite{gb_3d}.
\end{itemize}

В ходе параметрического моделирования вначале создается эскиз, с~которым впоследствии происходят изменения. В основе этого лежит математическая модель с подходящими параметрами, меняя значения которых, можно создать множество фигур. С помощью изменения параметров можно добиться необходимого вида модели.

Параметрическое моделирование появилось раньше полигонального как~совершенствование стандартной инженерной графики, способствует лучшему пониманию чертежей и зрительному восприятию разрабатываемых объектов.

Параметрическое моделирование применяется в промышленности\cite{sushentsov_2023}.
\subsection{Этапы 3D-моделирования}
На первом этапе создается пространственная геометрическая модель объекта, не учитывающая его физические характеристики. Производятся расчет размеров и построение формы предмета. Используются методы вращения, выдавливания, наращивания, полигонального моделирования.

На стадии создания текстуры объекта определяется, из каких материалов будет построен объект, разрабатывается его текстура. Именно в этот момент задается степень реалистичности создаваемой модели.

При выборе освещения возникают сложности, поскольку от указанных параметров зависит восприятие модели, насколько она будет правдоподобной. Указываются тон освещения, степень яркости, резкости, насыщенность теней.

На заключительной стадии 3D-моделирования осуществляется уточнение настроек отображения модели, в частности, добавление специальных эффектов вроде бликов, тумана и т.д. При наличии анимации корректируются ее параметры. Также определяются параметры визуализации (число кадров в~секунду, формат конечного видео). Если в результате получается двухмерное изображение, следует выбрать его формат и разрешение.

По окончании процесса 3D-моделирования в готовый материал можно включить спецэффекты с использованием программных средств, таких как Adobe Photoshop, Adobe Premier Pro, Adobe Illustrator и т.д Таким образом итоговый результат улучшается с применением различных технологий\cite{sky_3d}.
\subsection{Сущность фотограмметрии}
Фотограмметрия -- это технология получения точных пространственных данных об объектах на основе фотографических изображений. В контексте 3D моделирования фотограмметрия представляет собой процесс создания объемных цифровых моделей из множества двумерных снимков, сделанных с разных ракурсов. По сути, это математически обоснованный метод реконструкции трехмерного пространства на основе перекрывающихся изображений. 

В последние годы наблюдается растущий интерес и спрос на~методы фотограмметрии в реальном времени. Фотограмметрия в реальном времени -- это способность захватывать, обрабатывать и генерировать 3D-модели с~минимальными затратами времени, что позволяет визуализировать и~анализировать захваченную сцену. Эта развивающаяся область привлекла значительное внимание благодаря своим преимуществом в различных отраслях, виртуальную реальность, дополненную реальность и архитектуру.

Фотограмметрия в реальном времени находит применение в различных областях, таких как виртуальная реальность, дополненная реальность. Этот процесс обеспечивает интерактивный опыт, когда 3D-модели можно визуализировать и управлять ими в режиме реального времени, повышая погружение и вовлеченность пользователей. Кроме того, фотограмметрия в~реальном времени находит применение в сценариях, требующих быстрого принятия решений или немедленного доступа к 3D-данным, таких как архитектурная визуализация, мониторинг строительства и реагирование на~стихийные бедствия.

Ключевое отличие фотограмметрии от традиционных методов 3D моделирования заключается в подходе к созданию модели. Если при~классическом моделировании специалист вручную выстраивает геометрию с нуля, то фотограмметрия автоматически извлекает геометрические данные из реальных объектов через их фотографии\cite{incredibleart}.

Традиционные рабочие процессы фотограмметрии обычно включают этап постобработки, когда изображения обрабатываются в автономном режиме для реконструкции 3D-моделей. Однако последние достижения в~аппаратных возможностях и программных алгоритмах позволили разработать решения для фотограмметрии в реальном времени и на месте. Эти подходы предлагают множество преимуществ, включая немедленную обратную связь, более быструю обработку данных и возможность захвата и анализа 3D-информации непосредственно на месте. В этом разделе исследуется концепция фотограмметрии в реальном времени и на месте\cite{skillbox_fotogrammetriya}.

Принцип работы фотограмметрии основан на триангуляции -- процессе определения координат точек в трехмерном пространстве путем построения воображаемых треугольников от камеры к объекту. Когда одна и та же точка видна на двух или более снимках, сделанных с разных позиций, программа может вычислить её точное расположение в трехмерном пространстве.

Процесс включает захват и обработку 3D-данных непосредственно в~интересующем месте, что устраняет необходимость в обширной автономной обработке. Этот подход выигрывает от портативного компьютерного оборудования и специализированного программного обеспечения, которое может работать на месте, обеспечивая быстрое получение и реконструкцию данных. Фотограмметрия на месте особенно ценна в ситуациях, когда критически важен немедленный доступ к точной 3D-информации, например, при археологических раскопках, осмотре места преступления или мониторинге строительных площадок\cite{incredibleart}.
\subsection{Преимущества и недостатки фотограмметрии}
Бурное развитие потребительской электроники, в частности, повышение разрешающей способности и оптического качества камер в смартфонах и планшетах, существенно снизило порог входа в фотограмметрию. Это позволяет привлекать к процессу сбора исходных данных не только специалистов с профессиональным оборудованием, но и широкий круг неподготовленных пользователей, что способствует популяризации метода.

Несомненное преимущество фотограмметрии -- это экономия времени по~сравнению с обычным моделированием, что особенно чувствуется при~работе с масштабными и сложными по геометрии сценами или объектами, например, большими зданиями или целыми ландшафтными сценами. Благодаря программе, которая анализирует полученные снимки и создает трехмерную модель, можно избавиться от создания той же самой модели с~нуля и, получив полигональную модель, перейти к ее редактированию.

Помимо формы, технология способна определить по снимкам объекта его размер и положение в пространстве. 

Также преимуществом является возможность сделать свой проект по-настоящему уникальным. Особенно если прототипами игровых ассетов станут изделия ручной работы (игрушки, фигурки из пластилина и так далее), объекты архитектуры и культуры, которые могут подчеркнуть эпоху игры или национальный колорит сеттинга. Аудитория всегда проявляет интерес к~подобному контенту\cite{skillbox_fotogrammetriya}.

Однако у фотограмметрии не самая высокая точность. В этом отношении фотограмметрия уступает лазерному сканированию.

Получение профессиональных результатов сопряжено со значительными инвестициями. Это включает не только приобретение высококлассной фототехники и оборудования, но и лицензий на специализированное программное обеспечение для обработки данных (например, Agisoft Metashape, RealityCapture, Pix4D), стоимость которых может быть существенной\cite{sky_fotogrammetriya}.

Исходному скану в любом случае нужна доработка и оптимизация. Удаление лишних мешей, реконструкция не попавших в область скана участков (разрывы и слияния полигональной сетки), чистка геометрии, ремеш, создание новой UV-развёртки -- все эти операции требуют от пользователя дополнительного времени и более продвинутых навыков в программах для~3D-моделирования.

Текстурирование модели, полученной методом фотограмметрии, требует определённых художественных навыков. Когда объект сформирован в~3D-модель, на исходном материале могут образоваться различные артефакты текстуры в виде точек и растяжек. Инструмент клонирования не всегда помогает -- и в этом случае некоторые участки придётся дорисовывать вручную\cite{skillbox_fotogrammetriya}.

Несмотря на некоторые явные недостатки, фотограмметрия остаётся одним из лучших вариантов для проекта с реалистичной графикой, особенно при ограниченном бюджете. Фотограмметрию также можно использовать и в неигровых проектах -- например, для оформления визуализаций и~интерактивных онлайн-экспозиций. Она способна не просто продлить срок жизни, но и увековечить даже самый хрупкий предмет. А по мере развития технического прогресса фотограмметрия может стать обычным явлением в жизни\cite{incredibleart}.
\subsection{Оценка качества и анализ ошибок}
Оценка точности и надежности фотограмметрических реконструкций является неотъемлемым и решающим аспектом всего процесса. Очень важно количественно оценить качество реконструированных 3D-моделей и понять потенциальные ошибки и ограничения, связанные с фотограмметрическими рабочими процессами. Тщательная оценка качества помогает оценить пригодность полученных данных и убедиться, что реконструированные модели соответствуют требованиям точности\cite{sushentsov_2023}.

Оценка точности -- сравнение фотограмметрически реконструированных моделей с эталонными данными или другими надежными измерениями. Это можно сделать с помощью различных методов проверки, таких как прямые измерения, контрольные точки или сравнение с независимыми источниками данных. Количественно оценивая различия между эталонными данными и реконструированной моделью, можно определить точность фотограмметрических результатов.

На точность фотограмметрических результатов влияют несколько факторов. К ним относятся качество и разрешение входных изображений, точность калибровки камеры, алгоритмы выравнивания и сопоставления изображений, а также качество наземных опорных точек, используемых для географической привязки. Кроме того, такие факторы, как искажение объектива, атмосферные условия и геометрическая сложность сцены, могут внести ошибки, влияющие на точность конечной модели.

Анализ ошибок включает выявление и понимание источников ошибок в фотограмметрическом рабочем процессе. Он направлен на оценку надежности и ограничений реконструированных моделей. Анализируя различные источники ошибок, можно выявить потенциальные проблемы и~принять корректирующие меры для повышения общего качества процесса реконструкции. 

Ошибки в фотограмметрии могут возникать на разных этапах рабочего процесса, включая получение изображений, предварительную обработку изображений, извлечение признаков, сопоставление и создание плотного облака точек. Например, ошибки калибровки камеры или оценки искажения объектива могут распространяться на весь процесс реконструкции, влияя на точность конечной модели. Точно так же ошибки в алгоритмах обнаружения признаков и сопоставления могут привести к неточностям в~реконструированной геометрии. 

Важно проанализировать и количественно оценить эти ошибки, чтобы~понять ограничения рабочего процесса. Понимание источников ошибок позволяет исследователям и практикам разрабатывать стратегии для~уменьшения ошибок, уточнения алгоритмов и оптимизации процессов. Это обеспечивает производство более надежных и точных 3D-моделей.

Для количественной оценки фотограмметрических реконструкций используются различные метрики. Эти метрики оценивают геометрическую точность, полноту и уровень детализации реконструированных моделей. Обычно используемые метрики качества включают среднеквадратичную ошибку RMSE (root-mean-squared error), плотность облака точек, отклонение поверхности, оценки сопоставления признаков. Эти метрики обеспечивают объективные измерения\cite{sushentsov_2023}.

Оценка качества и анализ ошибок -- это фундаментальные шаги в~фотограмметрии для обеспечения точности, надежности реконструированных 3D-моделей. Оценивая точность реконструированных моделей, анализируя источники ошибок и используя соответствующие показатели качества, пользователи могут улучшить качество фотограмметрических реконструкций и обрести уверенность в надежности результатов.

Наиболее продуктивной представляется стратегия синергетического использования обоих методов, при которой фотограмметрия обеспечивает создание исходной геометрии, а средства 3D-моделирования позволяют выполнить оптимизацию структуры сетки и адаптацию к техническим требованиям\cite{incredibleart}.
\newpage
\section{Принципы функционирования технологий дополненной и~виртуальной реальности}
\subsection{Введение и фундаментальные понятия}
Виртуальная и дополненная реальность сегодня являются двумя основными ветвями единого направления -- расширенной реальности. Виртуальная реальность полностью замещает физический мир искусственно созданной средой и стремится к абсолютному эффекту присутствия пользователя внутри этой среды. Дополненная реальность сохраняет реальный мир в качестве базовой основы и лишь накладывает на него цифровые объекты, информацию и взаимодействия, обогащая восприятие, но не подменяя его. Такое различие определяет не только технические принципы работы каждой технологии, но и их философское значение: первая ведёт к временной изоляции от реальности и создаёт совершенно новый онтологический опыт, вторая усиливает и расширяет существующую реальность, делая её более информативной и функциональной. Именно поэтому исследователи часто рассматривают VR и AR не как конкурирующие, а как взаимодополняющие технологии, которые в перспективе сольются в единый континуум.
\subsection{Исторический контекст развития}
Идеи обеих технологий зародились ещё в 1960-е годы. Первые попытки создания стереоскопических дисплеев и систем отслеживания движений головы относятся к работам Мортона Хейлига (Sensorama, 1962) и Айвена Сазерланда («Меч Дамокла», 1968) -- тяжёлому шлему с механическим подвеской. Термин «виртуальная реальность» окончательно оформился в~конце 1980-х благодаря Джарону Ланье и его компании VPL Research. Однако настоящий технологический прорыв произошёл в 2012–2016 годах, когда появились доступные потребительские устройства (Oculus Rift DK1, HTC Vive, Samsung Gear VR, Google Cardboard) и мобильные платформы получили встроенные системы трекинга и сенсоры глубины. Именно в этот период технологии вышли из узких научных лабораторий и военных полигонов в массовое пользование, став частью повседневной жизни миллионов людей по всему миру.
\subsection{Принципы функционирования виртуальной реальности}
Работа VR построена на создании устойчивого ощущения полного присутствия человека в искусственном мире. Это ощущение возникает только тогда, когда мозг перестаёт различать виртуальное и реальное, а для этого система должна одновременно решать несколько критически важных задач. Изображение должно быть стереоскопическим и охватывать практически всё поле зрения пользователя, создавая эффект «окна в другой мир». Движение головы и тела должно отслеживаться в реальном времени с миллиметровой и угловой точностью -- современные системы достигают погрешности менее 0,1°. Задержка от движения до обновления картинки не должна превышать 20 миллисекунд, а лучше 7–12 миллисекунд, иначе возникает киберболезнь. Звук обязан быть пространственным и реагировать на повороты головы в реальном времени. Современные автономные шлемы достигают этих параметров благодаря мощным мобильным процессорам, дисплеям с частотой 90–120 Гц, сложным алгоритмам асинхронной репроекции.
\subsection{Принципы функционирования дополненной реальности}
Дополненная реальность работает принципиально иначе. Здесь главная цель -- не изолировать человека от мира, а максимально точно и естественно встроить виртуальные объекты в реальное пространство так, чтобы они воспринимались как его органичная часть. Устройство непрерывно сканирует окружающую среду, строит её трёхмерную карту в реальном времени (технология VSLAM) и определяет своё положение внутри этой карты с сантиметровой точностью. На полученную карту накладываются цифровые модели таким образом, чтобы они оставались неподвижными относительно реальных поверхностей даже при резком движении пользователя. Освещение виртуальных объектов автоматически подстраивается под реальное освещение сцены, тени падают правильно, а реальные предметы могут заслонять виртуальные (эффект окклюзии). Именно поэтому современные AR-устройства всё чаще оснащаются лидарами, стереокамерами, инерциальными блоками и нейронными ускорителями, способными обрабатывать огромные массивы данных в доли секунды.
\subsection{Общие фундаментальные принципы обеих технологий}
Несмотря на различия, VR и AR опираются на одни и те же базовые принципы, сформулированные ещё в 1965 году Айвеном Сазерландом в~его исторической работе. Система должна обеспечивать глубокое сенсорное погружение, мгновенное и естественное взаимодействие и стимулировать воображение пользователя качественным содержанием. Погружение достигается многосенсорной достоверностью -- не только зрением и слухом, но и, в перспективе, тактильными, термическими и даже обонятельными ощущениями. Взаимодействие возможно только при абсолютной точности и~скорости трекинга. Воображение активируется смысловой насыщенностью контента и его эмоциональной глубиной. Ещё один общий принцип -- обязательность шести степеней свободы (6 DoF): три поступательных и три вращательных. Только при их наличии мозг воспринимает пространство как~реальное и позволяет пользователю естественно двигаться внутри него.
\subsection{Аппаратная и программная архитектура}
Аппаратная часть современных систем включает дисплеи сверхвысокого разрешения (до 8K на глаз в прототипах), инерциальные измерительные блоки, многообъективные камеры, лидары и специализированные нейронные процессоры. Программная часть решает задачи рендеринга в реальном времени, предсказания движения пользователя на десятки миллисекунд вперёд, компенсации задержек и одновременного построения карты окружающего пространства. Ключевое требование -- минимальная латентность всей цепочки от сенсора до фотона, поэтому разработчики постоянно повышают частоту обновления экранов до 120–144 Гц, улучшают алгоритмы трекинга и переносят наиболее тяжёлые вычисления на специализированные чипы.
\subsection{Примеры практического применения}
В медицине Виртуальная реальность активно используется для обучения хирургов: студенты проводят сложные операции в полностью безопасной среде, где можно бесконечно повторять каждое движение и менять анатомические варианты. Дополненная реальность помогает врачам во время реальных операций, накладывая на пациента трёхмерные модели органов, сосудов и~опухолей прямо в поле зрения хирурга, что снижает риск ошибок и~сокращает время вмешательства.

В образовании школьники и студенты с помощью VR-шлемов оказываются внутри древнего Рима, на поверхности Марса или внутри молекулы ДНК, что превращает абстрактные знания в непосредственный опыт. AR-приложения позволяют учителю проецировать на обычную парту интерактивные модели солнечной системы, исторических событий или~химических реакций, которые дети могут трогать, вращать и разбирать.

В промышленности и строительстве инженеры используют AR-очки для визуализации будущих зданий прямо на стройплощадке в натуральную величину ещё до начала работ. Техники обслуживают сложное оборудование, видя пошаговые инструкции, схемы и предупреждения прямо на реальных деталях, что резко снижает количество ошибок и время простоя.

В культуре и развлечениях виртуальная реальность дарит возможность посетить концерты любимых артистов, которые уже не выступают, прогуляться по музеям, закрытым для публики. Pokémon GO стал первым по-настоящему массовым примером дополненной реальности, превратив обычные городские улицы в глобальное игровое пространство и показав, как AR может вовлекать миллионы людей в активное движение и социальное взаимодействие.
\subsection{Перспективы и направления развития}
Будущее технологий лежит в постепенном стирании границ между VR и AR и создании единого континуума расширенной реальности. Устройства становятся легче (цель -- вес обычных очков), автономнее и внешне всё больше похожи на повседневную оптику. Разрешение растёт до уровня сетчатки глаза, угол обзора приближается к естественному человеческому (около 200°). Появляются тактильные костюмы, перчатки, термические и даже обонятельные интерфейсы. На горизонте 2030-х годов -- нейроинтерфейсы прямого воздействия на мозг, которые позволят обойти традиционные органы чувств. Уже сегодня формируются первые персистентные метавселенные -- общие цифровые пространства, где миллионы людей одновременно живут, работают, учатся и творят.

Принципы функционирования виртуальной и дополненной реальности сводятся к созданию убедительной иллюзии присутствия через точное, быстрое и многосенсорное воспроизведение или дополнение реального мира. Главное различие заключается в степени радикальности этой иллюзии: полное замещение реальности в VR или её естественное расширение в AR. Сегодня эти технологии уже радикально меняют медицину, образование, промышленность, культуру и военное дело, а в ближайшие десятилетия они сольются в единую ткань человеческой деятельности, окончательно стерев границы между физическим и цифровым и создав новую среду существования человека.
\newpage
\section{Среды и инструменты для~разработки контента в~AR/VR}
Рассмотрим инструменты для создания контента в дополненной и~виртуальной реальности. Эти технологии стали новыми платформами для~выражения идей, где AR накладывает цифровые объекты на наш мир через экран смартфона или очков, а VR полностью погружает пользователя в~смоделированную среду с помощью шлема. Ядром процесса разработки для~обеих технологий являются комплексные среды разработки, самой известной из которых по праву считается Unity.

Unity представляет собой универсальную кроссплатформенную среду, построенную на принципах доступности и гибкости. Её интерфейс -- это настраиваемая рабочая область с ключевыми окнами: для визуального редактирования сцены (Scene), предварительного просмотра результата (Game), иерархии объектов (Hierarchy) и управления их свойствами (Inspector). Разработка в Unity начинается с визуальной сборки: 3D-модели, источники света, камеры и интерфейсы перетаскиваются в сцену и настраиваются. Вся интерактивная логика, от простой анимации до сложной игровой механики, создаётся путём написания скриптов на языке программирования C\#. Эти скрипты присоединяются к объектам как компоненты, что делает архитектуру проекта модульной и понятной. Важнейшей особенностью Unity является её адаптивность: разработчик может выбрать конвейер рендеринга, подходящий именно для его задачи. Например, Universal Render Pipeline оптимизирован для производительности в мобильных AR/VR-проектах, а High Definition Render Pipeline открывает путь к кинематографичному качеству графики. Неотъемлемой частью экосистемы является Unity Asset Store -- обширная библиотека готовых 3D-моделей, текстур, звуков и даже целых программных модулей, что значительно ускоряет процесс создания прототипа или полноценного приложения. Прямая интеграция с основными платформами через фреймворк OpenXR позволяет разработчику, создав проект один раз, легко адаптировать его для таких систем, как ARKit на iOS, ARCore на~Android или Oculus VR, делая Unity идеальным инструментом для быстрой и эффективной кроссплатформенной разработки.

Другой гигант в области создания интерактивного контента -- Unreal Engine. Эта среда сфокусирована на достижении высочайшего визуального качества и предоставляет художникам и дизайнерам мощные инструменты прямо «из коробки». Интерфейс Unreal Engine -- это профессиональная рабочая станция. Его отличают революционные графические системы, такие как Nanite, позволяющая импортировать в реальном времени невероятно детализированные 3D-модели из киноиндустрии, и Lumen, обеспечивающая динамическое реалистичное освещение, которое меняется в зависимости от~действий в сцене. Для создания интерактивной логики в~Unreal Engine активно используется система визуального программирования Blueprints. Она позволяет дизайнерам и техническим художникам создавать поведение объектов и целые игровые механики, соединяя узлы в графическом редакторе, что делает процесс прототипирования невероятно быстрым и~наглядным. Для~работы с материалами в движке есть мощный нодовый редактор, где художник, комбинируя различные текстуры и эффекты, может создать любую поверхность -- от потёртого металла до влажной кожи. Создание сложных визуальных эффектов, таких как огонь, дым или~магические частицы, доверено системе Niagara, которая предоставляет детальный контроль над поведением каждой частицы в симуляции. Благодаря такому набору инструментов Unreal Engine часто становится выбором для проектов, где критически важны фотореализм, сложная графика и~кинематографичность, будь то~высокобюджетные игры, профессиональные симуляторы или~инновационные VR-опыты.

Однако важно понимать, что сами движки редко используются для~первичного создания детализированного контента. Они выступают в~роли центрального хаба, куда стекаются и где оживают ресурсы, созданные в специализированных программах. Так, сложные 3D-модели и персонажи создаются в пакетах для моделирования и скульптинга, таких как Blender, Maya или ZBrush. Текстуры и сложные материалы для этих моделей разрабатываются в Substance Painter или Adobe Photoshop. Анимацию и движение зачастую настраивают в Maya или Blender, а дизайн интерфейсов и 2D-графику рисуют в Figma или Illustrator. Весь этот контент экспортируется в стандартных форматах и импортируется в среду разработки, где уже программируется его логика, настраивается взаимодействие и финальная сборка.

Таким образом, разработка для AR/VR -- это всегда синтез искусства и~инженерии, слаженная работа целого набора инструментов. Unity и Unreal Engine являются мощными универсальными средами, которые интегрируют в~себе работу всех остальных программ, предоставляя не только графический рендеринг, но и системы физики, звука, ввода и, что критически важно, прямые инструменты для работы с оборудованием AR/VR. Грамотное использование этих сред в связке со специализированным софтом -- это и есть ключ к созданию убедительного, стабильного и захватывающего цифрового опыта.
\newpage
\section{Применение AR и VR в~гуманитарных исследованиях и~образовании}
Виртуальная реальность (VR) изначально создавалась с целью полного погружения в искусственный мир, который передаётся человеку через его ощущения: зрение, слух, обоняние, осязание и другие. Поэтому сперва технологию опробовали в видеоиграх, где используются устройства ввода-вывода для глубокого погружения. 

Дополненная реальность (AR) изначально создавалась для дополнения реального мира виртуальными объектами с целью улучшения восприятия информации. В отличие от VR, где пользователь «перемещается» в другое место, AR накладывает цифровые элементы на реальное окружение, дополняя его. 

Таким образом, основная цель VR -- создание цифрового мира, максимально похожего на реальный, а AR -- расширение реального мира виртуальными объектами.

В последнее время цифровые технологии развиваются и в науке.
\subsection{Основные области применения в гуманитарных исследованиях}
VR/AR применяют для реконструкции архитектурных комплексов и~исторических пространств (панорамные VR-модели дворцов, виртуальные раскопки), что позволяет исследователям и широкой публике «посетить» потерянные или ограниченно доступные объекты. Такие реконструкции помогают тестировать гипотезы о планировке, интерьерах и использовании пространства. 

Виртуальные репродукции произведений, VR-экскурсии по историческим выставкам дают возможность детально изучать композицию, технику, контекст. Это особенно полезно для сравнительного анализа и для занятий по истории искусства. 

AR применяется для наложения метаданных, переводов и аннотаций прямо на оцифрованные тексты или рукописи; VR -- для создания «виртуальных архивов», где можно моделировать контекст хранения документов. Это упрощает междисциплинарный доступ к материалам. 

VR используется для моделирования социальных ситуаций (например, для изучения межгрупповых взаимодействий), где контролируемая имитация событий даёт преимущества перед традиционными методами. При этом важно учитывать этические аспекты подобных экспериментов. 
\subsection{Применение в образовании (гуманитарные дисциплины)}
AR/VR поддерживают конструктивистские и экспериментальные подходы: студенты не просто получают информацию, а взаимодействуют с~реконструированными средами (например, «прогулки» по древнему городу), что повышает вовлечённость и усвоение материала. Обзоры указывают на~улучшение мотивации и в ряде случаев -- учебных результатов. 

Музеи и проекты культурного наследия используют AR-метки и VR-туру для расширения доступа к коллекциям и создания интерактивных экспозиций. Это способствует популяризации науки и расширяет аудиторию. 

В условиях гибридного обучения VR-среды дают возможность создать единый учебный опыт для студентов, находящихся в разных местах, и~симулировать лабораторно-полевые занятия для дистанционных слушателей.

Преимущества AR/VR в гуманитарных проектах:
\begin{enumerate}[leftmargin=0em, itemindent=3.75em, nosep, label=\arabic*)]
\item Иммерсивность и эмпатия. VR позволяет «пережить» исторические или культурные ситуации, что может усиливать эмоциональное понимание материалов;
\item Визуализация сложных данных -- археологические слои, 3D-модели артефактов, пространственные реконструкции;
\item Доступ и репликация опыта -- объекты, недоступные по географии или~сохранности, становятся исследуемыми и демонстрируемыми. 
\end{enumerate}

Ограничения и проблемы внедрения:
\begin{enumerate}[leftmargin=0em, itemindent=3.75em, nosep, label=\arabic*)]
\item Технические и инфраструктурные барьеры -- стоимость оборудования, потребность в мощных компьютерах/серверах, проблемы масштабируемости;
\item Педагогические вызовы -- недостаток методик и готовых сценариев интеграции в учебные планы; требуется переквалификация преподавателей;
\item Этические вопросы -- воздействие иммерсивного опыта на психику, «манипуляция эмпатией», проблемы согласия при участии в экспериментах, конфиденциальность; требуется строгая этическая оценка исследований в VR;
\item Доступность и инклюзия -- риски усиления цифрового неравенства, проблемы для людей с ограничениями здоровья (кинетоз, проблемы зрения). 
\end{enumerate}

Практические примеры и кейсы:
\begin{enumerate}[leftmargin=0em, itemindent=3.75em, nosep, label=\arabic*)]
\item ERASMUS+ проект Virtual Worlds -- международный проект по~использованию VR в образовании для археологии, древних исследований и~истории искусства (пример сотрудничества академических центров и~музеев);
\item Панорамная VR-реконструкция Потала (китайские проекты культурного наследия) -- использование VR для массовой презентации архитектурных памятников;
\item Учебные курсы по истории искусства с VR -- кейсы университетских курсов, где VR-сессии интегрированы в практические занятия по анализу произведений.
\end{enumerate}
\subsection{Методические рекомендации для исследователей и преподавателей}
Определяйте учебные цели, а не технологию -- выбирайте AR/VR только если иммерсивность даёт добавленную ценность по сравнению с~традиционными методами. 

Проектируйте сценарии с учётом этики -- получайте информированное согласие, учитывайте психологические риски, обеспечивайте возможность выхода из сессии. 

Планируйте доступность -- альтернативные форматы для тех, кому недоступен VR (видео, 3D-модели на экране, описательные материалы). 

Оценивайте эффект количественно и качественно -- комбинируйте тесты знаний, анкетирование мотивации и поведенческие метрики внутри VR. 

AR и VR открывают значимые возможности для гуманитарных наук и образования: они помогают моделировать исчезнувшие пространства, углублять визуальный анализ, повышать вовлечённость студентов и~расширять доступ к культурным ресурсам. Вместе с тем успешное и этичное внедрение требует продуманной методологии, инвестиций в инфраструктуру и~подготовки преподавателей. Будущее направления -- интеграция AR/VR с цифровыми архивами, «оцифрованные двойники» утерянных объектов, стандартизованные педагогические сценарии для гуманитарной дисциплины.
\newpage
\section{Тенденции развития иммерсивных технологий и их влияние на современную культуру}
\subsection{Тенденции развития}
В последние годы иммерсивные технологии -- такие как виртуальная реальность (VR), дополненная реальность (AR) и смешанная реальность (MR) -- стремительно развиваются и становятся всё более распространёнными в~разных сферах. С 2014 по 2024 год число публикаций о применении VR, AR и MR в области культурного наследия существенно выросло, что показывает рост интереса научного сообщества к этим технологиям\cite{koh_impact_2024}. 

Музеи, галереи и культурные учреждения всё активнее интегрируют иммерсивные форматы: виртуальные туры, 360°-видео, MR-инсталляции. 

Также развивается направление «социальной виртуальной реальности» и цифровых культурных платформ, где пользователи из разных стран могут взаимодействовать, обсуждать искусство, посещать виртуальные выставки и~путешествовать по виртуальным репликам реальных мест.
\subsection{Влияние на современную культуру}
Виртуальные музеи позволяют посетителям удалённо исследовать экспонаты и взаимодействовать с ними. Это делает культурное наследие более инклюзивным и актуальным для различных аудиторий.  

Благодаря иммерсивным технологиям появились формы художественного выражения и творчества. Художники, режиссёры, курирующие виртуальные выставки, создают такие мультимедийные и интерактивные проекты, которые невозможно реализовать в традиционных условиях. 

Такие технологии влияют на то, как мы воспринимаем историю, наследие, культуру: благодаря виртуальным реконструкциям, аудио- и~визуальным эффектам, люди могут окунуться в прошлое, прочувствовать эпоху, увидеть объекты искусства и археологические памятники так, как будто они присутствуют лично. Это увеличивает эмоциональную и образовательную ценность культурных мероприятий\cite{ng_cultural_2025}.
\subsection{Риски и критика иммерсивных технологий}
Несмотря на широкий потенциал VR, AR и MR, развитие иммерсивных технологий сопровождается рядом рисков, которые важно учитывать.

Погружение в виртуальные среды делает цифровой опыт более привлекательным, чем взаимодействие в реальном мире. Это может усиливать разрыв между людьми, увеличивать зависимость от технологий и снижать количество живого общения.

Также критики указывают на опасность искажения реальности и~зависимости от виртуального опыта. Если человек регулярно предпочитает виртуальные миры реальному, это может привести к уходу от социальных и культурных задач, а также к формированию нереалистичных ожиданий от~жизни\cite{slater_enhancing_2016}.

Наконец, специалисты обсуждают проблему доступности: создание качественных иммерсивных проектов требует высоких затрат, а дорогое оборудование делает технологию недоступной для многих культурных организаций и пользователей.

Иммерсивные технологии -- VR, AR и MR -- активно развиваются и~превращаются в важный инструмент культуры и образования. Они создают новые формы доступа к искусству, культурному наследию и историческому опыту, расширяют границы творчества и дают возможность людям по-новому воспринимать мир. При разумном использовании они могут усилить культурную и образовательную значимость искусства, сделать её более доступной и эмоционально насыщенной.

Однако у иммерсивных технологий есть и недостатки: некоторые из~них могут привести к упрощению и обесцениванию искусства, превращая его в~развлекательный аттракцион. Кроме того, при некорректном или~избыточном применении иммерсивные технологии могут отрицательно сказаться на~ценностно-ориентационных воззрениях зрителя\cite{mironova_immersive_2023}.
\newpage
{
\sloppy
\nocite{*}
\cleardoublepage
\addcontentsline{toc}{section}{Список использованных источников}
\bibliographystyle{plainurl}
\bibliography{report10}
}
\end{document}